%% kapa: exam=1 round=26 year=2015 period=1 form=A subject=civil_law subject_ko=민법 seq=7 number=7
%% source: Q-Net 감정평가사 기출문제 (cst00309) artlSeq=5138001 file=extracted/1008/5138001/제26회_감정평가사 1차_기출/제26회_감정평가사 1차_기출/1교시_감평_A형.pdf
%% conversion: kordoc 4.24.51 -> tools/gen1.py (auto); review=proofread-sonnet; flags=none
\begin{question}{7}
착오로 인한 의사표시에 관한 설명으로 옳지 않은 것은? (다툼이 있으면 판례에 따름)
\begin{choices}
  \item 채무자의 동일성에 관한 물상보증인의 착오는 법률행위 내용의 중요부분에 관한 착오에 해당하지 않는다.
  \item 착오로 의사표시를 취소한 자는 이로 인하여 상대방에 대해 불법행위책임을 지지 않는다.
  \item 착오로 인하여 표의자가 경제적인 불이익을 입은 것이 아니라면 착오를 이유로 취소할 수 없다.
  \item 특별한 사정이 없는 한 목적물의 시가에 관한 착오를 이유로 매매계약을 취소할 수 없다.
  \item 동기의 착오가 상대방에 의해 유발된 경우 동기가 표시되지 않더라도 의사표시의 취소 사유인 착오에 해당할 수 있다.
\end{choices}
\end{question}
\answer{1}
